% --- Archon physics formalization source begin ---
% archon:physics
% archon:covers PhyXMiniProblems/problem_phyx_mini_0169.lean
% archon:source-report reports/phyx_mini/problem_phyx_mini_0169.source.json

\chapter{Physics problem phyx\_mini\_0169}
\label{ch:PhyXMiniProblems_problem_phyx_mini_0169}

\paragraph{Problem source.}
\#\# Physical scenario

Small speakers \textbackslash{}( A \textbackslash{}) and \textbackslash{}( B \textbackslash{}) are driven in phase at \textbackslash{}( 725\textbackslash{},\textbackslash{}text\{Hz\} \textbackslash{}) by the same audio oscillator. Both speakers start out \textbackslash{}( 4.50\textbackslash{},\textbackslash{}text\{m\} \textbackslash{}) from the listener, but speaker \textbackslash{}( A \textbackslash{}) is slowly moved away (\textbackslash{}textbf\{figure\}).

\#\# Question

At what distance \textbackslash{}( d \textbackslash{}) will the sound from the speakers first produce destructive interference at the listener’s location?

\#\# Answer choices

- A: A: 0.237m

- B: B: 0.336m

- C: C: 0.569m

- D: D: 0.414m

\#\# Figure caption (auxiliary; use the image as primary evidence)

The image is a diagram with several components:

1. **Points and Labels:**
   - Point **A** on the left, represented by a solid circle.
   - Point **B** to the right of A, also shown as a solid circle.
   - The distance between A and B is labeled as \textbackslash{}( d \textbackslash{}).

2. **Line and Measurements:**
   - A horizontal arrow extends from point B to the right.
   - The line from B to an illustrated face measures **4.50 m**.

3. **Illustrated Face:**
   - On the far right of the diagram, there is a sketch of a face with arrows pointing away from it, possibly indicating a viewpoint or direction.

The diagram seems to represent distances or directions between objects and a viewer or perspective.

\#\# Dataset metadata

Resonance and Harmonics; Physical Model Grounding Reasoning, Spatial Relation Reasoning

\paragraph{Recorded answer/context.}
A: A: 0.237m

\paragraph{Figure/image path.}
/root/proposal\_for\_physic/science-mango/phyx\_mini\_run/phyx\_data/test\_image/169.png

\paragraph{Formalization target.}
create a compiling Lean file with sorry bodies at `PhyXMiniProblems/problem\_phyx\_mini\_0169.lean`. The Lean declarations must preserve the physical quantities, dimensions or dimensional roles, figure labels, governing-law hypotheses, and final relation expressed by this problem.
Use Mathlib/Physlib names found through LeanExplore where available. If a physics API is missing, introduce faithful local abstractions rather than scalar placeholder aliases.

\begin{theorem}[Physics formalization target]
\label{thm:physics:phyx_mini_0169:target}
\lean{PhyXMiniProblems.ProblemPhyXMini0169.problem_phyx_mini_0169}
\uses{def:physics:phyx-mini-0169:phyxminiproblems-problemphyxmini0169-lengthinmeters, def:physics:phyx-mini-0169:phyxminiproblems-problemphyxmini0169-twospeakerinterferencesetup, def:physics:phyx-mini-0169:phyxminiproblems-problemphyxmini0169-driveninphasebysameoscillator, def:physics:phyx-mini-0169:phyxminiproblems-problemphyxmini0169-hasdepictedcollineargeometry, def:physics:phyx-mini-0169:phyxminiproblems-problemphyxmini0169-hasstatedreadouts, def:physics:phyx-mini-0169:phyxminiproblems-problemphyxmini0169-usesstandardambientairsoundspeed, def:physics:phyx-mini-0169:phyxminiproblems-problemphyxmini0169-hasphysicalparameters, def:physics:phyx-mini-0169:phyxminiproblems-problemphyxmini0169-satisfiesacousticinterferencelaws, def:physics:phyx-mini-0169:phyxminiproblems-problemphyxmini0169-isfirstdestructivedisplacement, def:physics:phyx-mini-0169:phyxminiproblems-problemphyxmini0169-answerchoice, def:physics:phyx-mini-0169:phyxminiproblems-problemphyxmini0169-matchesanswertonearestmillimeter}
The assigned autoformalize agent should translate this physics problem into Lean declarations in the covered file, with theorem and lemma proof bodies written as `by sorry`.
\end{theorem}
\begin{proof}
This is an autoformalization task, not a proof task. Produce statements that can later be proved by the physics prover without weakening the physical model.
\end{proof}

% --- Archon declaration topology begin ---
\section{Declaration topology}
\begin{definition}[Length Quantity]
  \label{def:physics:phyx-mini-0169:phyxminiproblems-problemphyxmini0169-lengthquantity}
  \lean{PhyXMiniProblems.ProblemPhyXMini0169.LengthQuantity}
  A physical quantity carrying the dimension of length
\end{definition}
\begin{proof}
  This is the declared model definition of Length Quantity.
\end{proof}
\begin{definition}[Frequency Quantity]
  \label{def:physics:phyx-mini-0169:phyxminiproblems-problemphyxmini0169-frequencyquantity}
  \lean{PhyXMiniProblems.ProblemPhyXMini0169.FrequencyQuantity}
  A physical frequency, carrying the inverse-time dimension
\end{definition}
\begin{proof}
  This is the declared model definition of Frequency Quantity.
\end{proof}
\begin{definition}[Speed Quantity]
  \label{def:physics:phyx-mini-0169:phyxminiproblems-problemphyxmini0169-speedquantity}
  \lean{PhyXMiniProblems.ProblemPhyXMini0169.SpeedQuantity}
  A physical propagation speed
\end{definition}
\begin{proof}
  This is the declared model definition of Speed Quantity.
\end{proof}
\begin{definition}[meters Value]
  \label{def:physics:phyx-mini-0169:phyxminiproblems-problemphyxmini0169-metersvalue}
  \lean{PhyXMiniProblems.ProblemPhyXMini0169.metersValue}
  \uses{def:physics:phyx-mini-0169:phyxminiproblems-problemphyxmini0169-lengthquantity}
  Read a physical length in metres
\end{definition}
\begin{proof}
  This is the declared model definition of meters Value.
\end{proof}
\begin{definition}[hertz Value]
  \label{def:physics:phyx-mini-0169:phyxminiproblems-problemphyxmini0169-hertzvalue}
  \lean{PhyXMiniProblems.ProblemPhyXMini0169.hertzValue}
  \uses{def:physics:phyx-mini-0169:phyxminiproblems-problemphyxmini0169-frequencyquantity}
  Read a physical frequency in hertz
\end{definition}
\begin{proof}
  This is the declared model definition of hertz Value.
\end{proof}
\begin{definition}[meters Per Second Value]
  \label{def:physics:phyx-mini-0169:phyxminiproblems-problemphyxmini0169-meterspersecondvalue}
  \lean{PhyXMiniProblems.ProblemPhyXMini0169.metersPerSecondValue}
  \uses{def:physics:phyx-mini-0169:phyxminiproblems-problemphyxmini0169-speedquantity}
  Read a physical speed in metres per second
\end{definition}
\begin{proof}
  This is the declared model definition of meters Per Second Value.
\end{proof}
\begin{definition}[length In Meters]
  \label{def:physics:phyx-mini-0169:phyxminiproblems-problemphyxmini0169-lengthinmeters}
  \lean{PhyXMiniProblems.ProblemPhyXMini0169.lengthInMeters}
  \uses{def:physics:phyx-mini-0169:phyxminiproblems-problemphyxmini0169-lengthquantity}
  Construct the physical length whose SI readout is the supplied number of metres
\end{definition}
\begin{proof}
  This is the declared model definition of length In Meters.
\end{proof}
\begin{definition}[Speaker Label]
  \label{def:physics:phyx-mini-0169:phyxminiproblems-problemphyxmini0169-speakerlabel}
  \lean{PhyXMiniProblems.ProblemPhyXMini0169.SpeakerLabel}
  Labels of the two acoustic sources
\end{definition}
\begin{proof}
  This is the declared model definition of Speaker Label.
\end{proof}
\begin{definition}[Figure Point]
  \label{def:physics:phyx-mini-0169:phyxminiproblems-problemphyxmini0169-figurepoint}
  \lean{PhyXMiniProblems.ProblemPhyXMini0169.FigurePoint}
  Point labels visible in the primary figure
\end{definition}
\begin{proof}
  This is the declared model definition of Figure Point.
\end{proof}
\begin{definition}[figure Point]
  \label{def:physics:phyx-mini-0169:phyxminiproblems-problemphyxmini0169-speakerlabel-figurepoint}
  \lean{PhyXMiniProblems.ProblemPhyXMini0169.SpeakerLabel.figurePoint}
  \uses{def:physics:phyx-mini-0169:phyxminiproblems-problemphyxmini0169-speakerlabel, def:physics:phyx-mini-0169:phyxminiproblems-problemphyxmini0169-figurepoint}
  The figure point occupied by a labeled speaker
\end{definition}
\begin{proof}
  This is the declared model definition of figure Point.
\end{proof}
\begin{definition}[Speaker Motion Regime]
  \label{def:physics:phyx-mini-0169:phyxminiproblems-problemphyxmini0169-speakermotionregime}
  \lean{PhyXMiniProblems.ProblemPhyXMini0169.SpeakerMotionRegime}
  The slow-motion regime described in the prose
\end{definition}
\begin{proof}
  This is the declared model definition of Speaker Motion Regime.
\end{proof}
\begin{definition}[Acoustic Speaker]
  \label{def:physics:phyx-mini-0169:phyxminiproblems-problemphyxmini0169-acousticspeaker}
  \lean{PhyXMiniProblems.ProblemPhyXMini0169.AcousticSpeaker}
  \uses{def:physics:phyx-mini-0169:phyxminiproblems-problemphyxmini0169-frequencyquantity}
  A monochromatic speaker at the common oscillator reference time
\end{definition}
\begin{proof}
  This is the declared model definition of Acoustic Speaker.
\end{proof}
\begin{definition}[Two Speaker Interference Setup]
  \label{def:physics:phyx-mini-0169:phyxminiproblems-problemphyxmini0169-twospeakerinterferencesetup}
  \lean{PhyXMiniProblems.ProblemPhyXMini0169.TwoSpeakerInterferenceSetup}
  \uses{def:physics:phyx-mini-0169:phyxminiproblems-problemphyxmini0169-lengthquantity, def:physics:phyx-mini-0169:phyxminiproblems-problemphyxmini0169-frequencyquantity, def:physics:phyx-mini-0169:phyxminiproblems-problemphyxmini0169-speedquantity, def:physics:phyx-mini-0169:phyxminiproblems-problemphyxmini0169-speakerlabel, def:physics:phyx-mini-0169:phyxminiproblems-problemphyxmini0169-figurepoint, def:physics:phyx-mini-0169:phyxminiproblems-problemphyxmini0169-speakermotionregime, def:physics:phyx-mini-0169:phyxminiproblems-problemphyxmini0169-acousticspeaker}
  The physical apparatus and its observables. The argument of figurePositionMetersAt, rayPathLengthAt, and destructiveInterferenceAt is the nonnegative physical displacement of speaker A from its initial position. In the primary image that same quantity is the speaker separation labelled d
\end{definition}
\begin{proof}
  This is the declared model definition of Two Speaker Interference Setup.
\end{proof}
\begin{definition}[Driven In Phase By Same Oscillator]
  \label{def:physics:phyx-mini-0169:phyxminiproblems-problemphyxmini0169-driveninphasebysameoscillator}
  \lean{PhyXMiniProblems.ProblemPhyXMini0169.DrivenInPhaseBySameOscillator}
  \uses{def:physics:phyx-mini-0169:phyxminiproblems-problemphyxmini0169-twospeakerinterferencesetup}
  Both speakers are driven at the oscillator frequency and with the same phase
\end{definition}
\begin{proof}
  This is the declared model definition of Driven In Phase By Same Oscillator.
\end{proof}
\begin{definition}[speaker Separation Meters At]
  \label{def:physics:phyx-mini-0169:phyxminiproblems-problemphyxmini0169-speakerseparationmetersat}
  \lean{PhyXMiniProblems.ProblemPhyXMini0169.speakerSeparationMetersAt}
  \uses{def:physics:phyx-mini-0169:phyxminiproblems-problemphyxmini0169-lengthquantity, def:physics:phyx-mini-0169:phyxminiproblems-problemphyxmini0169-twospeakerinterferencesetup}
  Speaker separation in the signed metre chart of the primary figure
\end{definition}
\begin{proof}
  This is the declared model definition of speaker Separation Meters At.
\end{proof}
\begin{definition}[Has Depicted Collinear Geometry]
  \label{def:physics:phyx-mini-0169:phyxminiproblems-problemphyxmini0169-hasdepictedcollineargeometry}
  \lean{PhyXMiniProblems.ProblemPhyXMini0169.HasDepictedCollinearGeometry}
  \uses{def:physics:phyx-mini-0169:phyxminiproblems-problemphyxmini0169-lengthquantity, def:physics:phyx-mini-0169:phyxminiproblems-problemphyxmini0169-metersvalue, def:physics:phyx-mini-0169:phyxminiproblems-problemphyxmini0169-speakerlabel, def:physics:phyx-mini-0169:phyxminiproblems-problemphyxmini0169-twospeakerinterferencesetup, def:physics:phyx-mini-0169:phyxminiproblems-problemphyxmini0169-speakerseparationmetersat}
  The collinear layout read from the primary raster image. The listener is at coordinate zero, B stays 4.50 m to the left once the numerical readout is supplied, and moving A away by d puts it a further d to the left. Consequently the labeled A--B separation is d, and each ray length is the absolute source-to-listener coordinate difference. At displacement zero this also states that both speakers start at the same distance from the listener
\end{definition}
\begin{proof}
  This is the declared model definition of Has Depicted Collinear Geometry.
\end{proof}
\begin{definition}[Has Stated Readouts]
  \label{def:physics:phyx-mini-0169:phyxminiproblems-problemphyxmini0169-hasstatedreadouts}
  \lean{PhyXMiniProblems.ProblemPhyXMini0169.HasStatedReadouts}
  \uses{def:physics:phyx-mini-0169:phyxminiproblems-problemphyxmini0169-metersvalue, def:physics:phyx-mini-0169:phyxminiproblems-problemphyxmini0169-hertzvalue, def:physics:phyx-mini-0169:phyxminiproblems-problemphyxmini0169-twospeakerinterferencesetup}
  Numerical readouts stated in the prose and figure. This predicate contains the 725 Hz oscillator frequency, the 4.50 m initial distance, and the quasistatic direction of motion, but no destructive-interference conclusion
\end{definition}
\begin{proof}
  This is the declared model definition of Has Stated Readouts.
\end{proof}
\begin{definition}[Uses Standard Ambient Air Sound Speed]
  \label{def:physics:phyx-mini-0169:phyxminiproblems-problemphyxmini0169-usesstandardambientairsoundspeed}
  \lean{PhyXMiniProblems.ProblemPhyXMini0169.UsesStandardAmbientAirSoundSpeed}
  \uses{def:physics:phyx-mini-0169:phyxminiproblems-problemphyxmini0169-meterspersecondvalue, def:physics:phyx-mini-0169:phyxminiproblems-problemphyxmini0169-twospeakerinterferencesetup}
  Ambient-air calibration used by the recorded answer. The source does not state temperature or sound speed; 343 m/s is the conventional room- temperature model which gives the listed 0.237 m after rounding
\end{definition}
\begin{proof}
  This is the declared model definition of Uses Standard Ambient Air Sound Speed.
\end{proof}
\begin{definition}[Has Physical Parameters]
  \label{def:physics:phyx-mini-0169:phyxminiproblems-problemphyxmini0169-hasphysicalparameters}
  \lean{PhyXMiniProblems.ProblemPhyXMini0169.HasPhysicalParameters}
  \uses{def:physics:phyx-mini-0169:phyxminiproblems-problemphyxmini0169-lengthquantity, def:physics:phyx-mini-0169:phyxminiproblems-problemphyxmini0169-metersvalue, def:physics:phyx-mini-0169:phyxminiproblems-problemphyxmini0169-hertzvalue, def:physics:phyx-mini-0169:phyxminiproblems-problemphyxmini0169-meterspersecondvalue, def:physics:phyx-mini-0169:phyxminiproblems-problemphyxmini0169-speakerlabel, def:physics:phyx-mini-0169:phyxminiproblems-problemphyxmini0169-twospeakerinterferencesetup}
  Positivity conditions for the physical length, frequency, and speed parameters
\end{definition}
\begin{proof}
  This is the declared model definition of Has Physical Parameters.
\end{proof}
\begin{definition}[path Difference Meters]
  \label{def:physics:phyx-mini-0169:phyxminiproblems-problemphyxmini0169-pathdifferencemeters}
  \lean{PhyXMiniProblems.ProblemPhyXMini0169.pathDifferenceMeters}
  \uses{def:physics:phyx-mini-0169:phyxminiproblems-problemphyxmini0169-lengthquantity, def:physics:phyx-mini-0169:phyxminiproblems-problemphyxmini0169-metersvalue, def:physics:phyx-mini-0169:phyxminiproblems-problemphyxmini0169-twospeakerinterferencesetup}
  Absolute difference of the two speaker-to-listener path lengths, in metres
\end{definition}
\begin{proof}
  This is the declared model definition of path Difference Meters.
\end{proof}
\begin{definition}[Satisfies Acoustic Interference Laws]
  \label{def:physics:phyx-mini-0169:phyxminiproblems-problemphyxmini0169-satisfiesacousticinterferencelaws}
  \lean{PhyXMiniProblems.ProblemPhyXMini0169.SatisfiesAcousticInterferenceLaws}
  \uses{def:physics:phyx-mini-0169:phyxminiproblems-problemphyxmini0169-lengthquantity, def:physics:phyx-mini-0169:phyxminiproblems-problemphyxmini0169-metersvalue, def:physics:phyx-mini-0169:phyxminiproblems-problemphyxmini0169-hertzvalue, def:physics:phyx-mini-0169:phyxminiproblems-problemphyxmini0169-meterspersecondvalue, def:physics:phyx-mini-0169:phyxminiproblems-problemphyxmini0169-twospeakerinterferencesetup, def:physics:phyx-mini-0169:phyxminiproblems-problemphyxmini0169-driveninphasebysameoscillator, def:physics:phyx-mini-0169:phyxminiproblems-problemphyxmini0169-pathdifferencemeters}
  The governing laws of the idealized two-source acoustic model: the nondispersive relation is speed = frequency * wavelength in SI; for coherent in-phase sources, a nonnegative displacement is destructive exactly when its path difference is an odd multiple of half a wavelength. Both clauses are uniform laws. In particular, neither names a first minimum, a numerical displacement, or an answer choice
\end{definition}
\begin{proof}
  This is the declared model definition of Satisfies Acoustic Interference Laws.
\end{proof}
\begin{definition}[Is First Destructive Displacement]
  \label{def:physics:phyx-mini-0169:phyxminiproblems-problemphyxmini0169-isfirstdestructivedisplacement}
  \lean{PhyXMiniProblems.ProblemPhyXMini0169.IsFirstDestructiveDisplacement}
  \uses{def:physics:phyx-mini-0169:phyxminiproblems-problemphyxmini0169-lengthquantity, def:physics:phyx-mini-0169:phyxminiproblems-problemphyxmini0169-metersvalue, def:physics:phyx-mini-0169:phyxminiproblems-problemphyxmini0169-twospeakerinterferencesetup}
  A displacement is the first destructive one when it is positive and destructive and no smaller positive physical displacement is destructive. This is a behavioral definition over the setup observable, not a numerical answer by definition
\end{definition}
\begin{proof}
  This is the declared model definition of Is First Destructive Displacement.
\end{proof}
\begin{definition}[Answer Choice]
  \label{def:physics:phyx-mini-0169:phyxminiproblems-problemphyxmini0169-answerchoice}
  \lean{PhyXMiniProblems.ProblemPhyXMini0169.AnswerChoice}
  Labels of the four distances displayed as answers
\end{definition}
\begin{proof}
  This is the declared model definition of Answer Choice.
\end{proof}
\begin{definition}[displayed Distance Meters]
  \label{def:physics:phyx-mini-0169:phyxminiproblems-problemphyxmini0169-displayeddistancemeters}
  \lean{PhyXMiniProblems.ProblemPhyXMini0169.displayedDistanceMeters}
  \uses{def:physics:phyx-mini-0169:phyxminiproblems-problemphyxmini0169-answerchoice}
  Metre readout printed beside an answer label
\end{definition}
\begin{proof}
  This is the declared model definition of displayed Distance Meters.
\end{proof}
\begin{definition}[Matches Answer To Nearest Millimeter]
  \label{def:physics:phyx-mini-0169:phyxminiproblems-problemphyxmini0169-matchesanswertonearestmillimeter}
  \lean{PhyXMiniProblems.ProblemPhyXMini0169.MatchesAnswerToNearestMillimeter}
  \uses{def:physics:phyx-mini-0169:phyxminiproblems-problemphyxmini0169-answerchoice, def:physics:phyx-mini-0169:phyxminiproblems-problemphyxmini0169-displayeddistancemeters}
  An exact physical distance matches a value printed to the nearest millimetre when the two metre readouts differ by at most half a millimetre
\end{definition}
\begin{proof}
  This is the declared model definition of Matches Answer To Nearest Millimeter.
\end{proof}
\begin{lemma}[meters Value length In Meters]
  \label{lem:physics:phyx-mini-0169:phyxminiproblems-problemphyxmini0169-metersvalue-lengthinmeters}
  \lean{PhyXMiniProblems.ProblemPhyXMini0169.metersValue_lengthInMeters}
  \uses{def:physics:phyx-mini-0169:phyxminiproblems-problemphyxmini0169-metersvalue, def:physics:phyx-mini-0169:phyxminiproblems-problemphyxmini0169-lengthinmeters}
  SI construction really has the requested metre readout
\end{lemma}
\begin{proof}
  The conclusion follows from the governing relations and figure readouts named in the statement of meters Value length In Meters.
\end{proof}
\begin{lemma}[path Difference Meters eq displacement]
  \label{lem:physics:phyx-mini-0169:phyxminiproblems-problemphyxmini0169-pathdifferencemeters-eq-displacement}
  \lean{PhyXMiniProblems.ProblemPhyXMini0169.pathDifferenceMeters_eq_displacement}
  \uses{def:physics:phyx-mini-0169:phyxminiproblems-problemphyxmini0169-lengthquantity, def:physics:phyx-mini-0169:phyxminiproblems-problemphyxmini0169-metersvalue, def:physics:phyx-mini-0169:phyxminiproblems-problemphyxmini0169-twospeakerinterferencesetup, def:physics:phyx-mini-0169:phyxminiproblems-problemphyxmini0169-hasdepictedcollineargeometry, def:physics:phyx-mini-0169:phyxminiproblems-problemphyxmini0169-pathdifferencemeters}
  In the depicted one-dimensional geometry, moving A away by d adds d to its listener path while the path from B remains fixed. Thus the path difference is exactly the displayed separation d
\end{lemma}
\begin{proof}
  The conclusion follows from the governing relations and figure readouts named in the statement of path Difference Meters eq displacement.
\end{proof}
\begin{lemma}[wavelength meters eq]
  \label{lem:physics:phyx-mini-0169:phyxminiproblems-problemphyxmini0169-wavelength-meters-eq}
  \lean{PhyXMiniProblems.ProblemPhyXMini0169.wavelength_meters_eq}
  \uses{def:physics:phyx-mini-0169:phyxminiproblems-problemphyxmini0169-metersvalue, def:physics:phyx-mini-0169:phyxminiproblems-problemphyxmini0169-twospeakerinterferencesetup, def:physics:phyx-mini-0169:phyxminiproblems-problemphyxmini0169-hasstatedreadouts, def:physics:phyx-mini-0169:phyxminiproblems-problemphyxmini0169-usesstandardambientairsoundspeed, def:physics:phyx-mini-0169:phyxminiproblems-problemphyxmini0169-hasphysicalparameters, def:physics:phyx-mini-0169:phyxminiproblems-problemphyxmini0169-satisfiesacousticinterferencelaws}
  The calibrated speed and stated frequency determine the wavelength to be 343 / 725 m by the acoustic dispersion relation
\end{lemma}
\begin{proof}
  The conclusion follows from the governing relations and figure readouts named in the statement of wavelength meters eq.
\end{proof}
\begin{lemma}[exact first destructive displacement]
  \label{lem:physics:phyx-mini-0169:phyxminiproblems-problemphyxmini0169-exact-first-destructive-displacement}
  \lean{PhyXMiniProblems.ProblemPhyXMini0169.exact_first_destructive_displacement}
  \uses{def:physics:phyx-mini-0169:phyxminiproblems-problemphyxmini0169-lengthinmeters, def:physics:phyx-mini-0169:phyxminiproblems-problemphyxmini0169-twospeakerinterferencesetup, def:physics:phyx-mini-0169:phyxminiproblems-problemphyxmini0169-driveninphasebysameoscillator, def:physics:phyx-mini-0169:phyxminiproblems-problemphyxmini0169-hasdepictedcollineargeometry, def:physics:phyx-mini-0169:phyxminiproblems-problemphyxmini0169-hasstatedreadouts, def:physics:phyx-mini-0169:phyxminiproblems-problemphyxmini0169-usesstandardambientairsoundspeed, def:physics:phyx-mini-0169:phyxminiproblems-problemphyxmini0169-hasphysicalparameters, def:physics:phyx-mini-0169:phyxminiproblems-problemphyxmini0169-satisfiesacousticinterferencelaws, def:physics:phyx-mini-0169:phyxminiproblems-problemphyxmini0169-isfirstdestructivedisplacement}
  The order-zero odd half-wavelength is 343 / 1450 m; the uniform interference criterion and the positivity hypotheses make it the first positive destructive displacement
\end{lemma}
\begin{proof}
  The conclusion follows from the governing relations and figure readouts named in the statement of exact first destructive displacement.
\end{proof}
% --- Archon declaration topology end ---
% --- Archon physics formalization source end ---
